\section{Method}
\label{sec:method}

We frame the method as a \emph{controlled diagnostic design} for partial observability, not as a search for the most powerful policy class. The key identification choice is therefore to hold the control policy class fixed while varying only the retained-state mechanism. All policy variants use the same finite-state-machine (FSM) controller, the same hand-written transition logic, the same action primitives, and the same observation interface; they differ only in how previously seen task-relevant information is retained after the goal becomes hidden. This fixed FSM is a deliberate controlled scaffold, not a methodological shortcut. If we changed controller capacity, exploration behavior, or action vocabularies together with memory, any success difference could always be attributed to increased controller capacity, leaving the memory question unanswered. By freezing those factors, we make retained state the only manipulated variable and can interpret the resulting memory gap as controlled evidence about \emph{what information must survive occlusion} rather than about which controller is more expressive \citep{Kaelbling1998}. Our claim is therefore intentionally narrower but cleaner: within a matched control backbone, does persistence alone repair the PO failure, and where does that repair stop working?

\subsection{Controlled Diagnostic Family}

\begin{itemize}
\item \task{NoMemory} uses the shared controller without retaining goal information after the initial observation.
\item \task{TextBuffer} uses the same controller and primitives, but retains the goal position in a lightweight flat buffer.
\item \task{StructMemory} uses the same controller and primitives, but retains the same core content in a key-value format with typed fields and event records (FSM state, retry count, rollback flag).
\end{itemize}

All three variants execute the same five-stage grasp sequence (approach $\to$ descend $\to$ grasp $\to$ lift/move $\to$ release) through the same fixed action primitives. This shared controller should not be read as assuming that real robots ought to use FSMs universally. It is the minimal matched scaffold that makes the diagnostic identifiable: if a flat cache and a structured store produce different outcomes under the same primitives, the difference is attributable to retained state; if they do not, then richer memory formatting is not the active ingredient in this benchmark family. The variant-to-question mapping is:
\begin{center}
\footnotesize
\setlength{\tabcolsep}{5pt}
\begin{tabular}{lll}
\hline
Variant & Question & Diagnostic target \\
\hline
\task{NoMemory} & Q1 & memory necessity \\
\task{TextBuffer} & Q2 & point-memory sufficiency \\
\task{StructMemory} & Q3 & beyond-persistence benefit \\
\hline
\end{tabular}
\end{center}
Equivalently, \task{NoMemory} asks whether retained state is needed at all, \task{TextBuffer} asks whether simple persistence is sufficient, and \task{StructMemory} asks whether richer organization improves over simple persistence. Because the retained-state mechanism is the only changing factor, contrasts among these variants directly implement the three diagnostic questions rather than merely comparing unrelated policies. Our retry ablations later probe the same Q3 axis from the control side by asking whether additional recovery logic helps, or instead consumes finite-horizon budget.

\subsection{Partial-Observability Protocol}

The PO wrapper masks task-relevant observations after the initial step and exposes the goal for one step every \texttt{flash\_every} steps (default: 20), controlled by \texttt{flash\_len} (default: 1). We define two diagnostic protocols:
\begin{itemize}
\item \textbf{Weak PO}: masks goal position only, so any success difference isolates the need to retain previously observed goal information.
\item \textbf{Strong PO}: masks both goal and object positions, so the diagnosis tests whether performance also depends on retained manipulated-object state rather than goal information alone.
\end{itemize}
Figure~\ref{fig:flash_every} in the appendix varies \texttt{flash\_every} across $\{1, 5, 10, 20, 50\}$ to control PO severity directly.

\subsection{Controlled Diagnostic Logic}

We formalize three progressively stronger diagnostic layers as a strict hierarchy:

\begin{definition}[Memory Necessity]
The \emph{memory gap} is
\[
\Delta_{\text{mem}}(\tau) = \max_{\pi \in \{\text{TB}, \text{SM}\}} S(\pi, \tau, \text{PO}) - S(\text{NM}, \tau, \text{PO}).
\]
Memory is \emph{necessary} iff $\Delta_{\text{mem}}(\tau) > 0$ with $p < 0.05$.
\end{definition}

Because \task{NoMemory}, \task{TextBuffer}, and \task{StructMemory} share the same FSM controller and fixed action primitives, while differing only in retained-state mechanism, $\Delta_{\text{mem}}(\tau)$ supports a controlled-ablation interpretation inside this benchmark family: it measures the success gain associated with adding retention to an otherwise unchanged controller.

\begin{definition}[Structure Differentiation]
The \emph{structure gap} is
\[
\Delta_{\text{struct}}(\tau) = |S(\text{TB}, \tau, \text{PO}) - S(\text{SM}, \tau, \text{PO})|.
\]
Structure \emph{differentiates} iff $\Delta_{\text{struct}}(\tau) > 0$ with significance; otherwise, the task requires only \emph{point memory}.
\end{definition}

\begin{definition}[Goal-Dependency Gradient]
\[
G(\tau) = S(\text{NM}, \tau, \text{FO}) - S(\text{NM}, \tau, \text{PO}).
\]
We use $G(\tau)$ as a \emph{diagnostic screening signal}: in our 9-task suite, $G(\tau) > 0.5$ marks \emph{high-dependency} tasks and $G(\tau) < 0.1$ marks \emph{low-dependency} tasks.
\end{definition}

\begin{proposition}[Diagnostic Layers]
Layer 1 tests whether hidden information matters ($\Delta_{\text{mem}}(\tau) > 0$). Layer 2 tests whether simple persistence is sufficient ($S(\text{TB}, \tau, \text{PO}) \approx S(\text{TB}, \tau, \text{FO})$). Layer 3 tests whether richer organization adds value ($\Delta_{\text{struct}}(\tau) > 0$). Each diagnostic layer subsumes the previous one.
\end{proposition}

Together, these diagnostic layers define a lightweight controlled workflow. We first estimate or measure the \task{NoMemory} FO-to-PO gap $G(\tau)$ as a screening signal for whether a memory comparison is worth running at all. We then interpret $\Delta_{\text{mem}}(\tau)$ as the gain associated with adding retained state to the shared controller, and $\Delta_{\text{struct}}(\tau)$ as the incremental gain associated with organizing that retained state more richly. The diagnostic outcome therefore localizes the dominant bottleneck to one of three places: observability itself, memory persistence, or memory structure.

\subsection{Goal-Relative Extension}

For contact-rich tasks, we introduce a \emph{Goal-Relative} extension: every goal-oriented stage reads \texttt{obs.goal\_pos} at every step. When the goal is masked (zero), the stage returns zero action. This prevents cached-goal leakage; otherwise, the shared FSM controller could bypass the PO mask via previously cached goal information.
